
\section{Background}
A robot may win inference time yet receive no usable result: its action queue can expire, or a moving target can leave the allowed pose tolerance. RTC studies delayed action execution \cite{rtc}; Armory studies heterogeneous robot policy serving \cite{armory}; Themis already links GPU allocation and fairness \cite{themis}. Our gap is a bounded comparison of usable service, private-value incentives and access policies, not the invention of deadline-aware scheduling.

\section{Motivation}
Charging for an expired result changes the incentive to participate. A nominal allocation therefore cannot stand in for usable access. The policy decision is how to distribute timely service and disclose the value cost of entrant priority and worker evaluation reserves.

\section{Application Scenario}
Independent small integrators and incumbent operators share a tabletop pick-and-place inference server. A coordinator commits to admission and payments; workers select tasks for reserved evaluation time. Verified telemetry fixes physical feasibility. Local fallback and protective control lie outside the auction; timely inference does not certify robot safety.

\section{Research Question}
\textbf{Shared question:} Can an access policy allocate usable robot inference while preserving truthful value reports? \textbf{Economics:} how do task value and entrant access trade off? \textbf{Computation:} which state constraints and reports change outcomes? \textbf{Behavior:} can users distinguish payment from timely completion? The final self-challenge tests the boundary between verifiable state and strategic reports.

\section{Method}
All jobs arrive at time zero at one nonpreemptive server. Operator $i$ privately values a timely result at $v_i\geq0$ and submits a sealed bid $b_i\geq0$. Public data specify service time, queued steps $h_i$, control period $\Delta_i$, pose tolerance $\epsilon_i$, speed bound $w_i$, network delay $\eta_i$ and guard time $g_i$. In milliseconds,
\[
d_i=\min\{h_i\Delta_i,1000\epsilon_i/w_i\}-\eta_i-g_i.
\]
For $w_i=0$, pose lifetime is infinite. Enumerate subsets, check earliest-deadline-first (EDF) completion against $d_i$ and horizon $T-r$, and maximize $\sum_i(b_i+ce_i)x_i$. Here $r$ reserves tail time and $e_i$ is verified entrant status. Winners pay $p_i=\max\{0,H_i^0-H_i^1-ce_i\}$, with best rival scores $H_i^0$ excluding $i$ and $H_i^1$ feasible with $i$. Losers pay zero; a fixed tie rule applies. Execution and settlement end the round.

\textbf{Game theory:} utility is $v_i$ for timely delivery minus payment. Fixed public feasibility gives dominant-strategy truthful values and a truthful Bayesian equilibrium under any private-value prior \cite{roughgarden,harsanyi}. \textbf{Social choice:} report task value, revenue, usable entrant access and reserved time separately; credits express policy priority, not measured social benefit. \textbf{Mechanism design:} compare capacity-only admission, EDF without payments, feasible admission, credits and reserves. Capacity-only admission is a diagnostic, not a physical-equilibrium benchmark. Exact enumeration is an offline oracle limited to 12 jobs.

\section{Results}
\begin{center}\small\setlength{\tabcolsep}{4pt}
\begin{tabular}{lcccc}\toprule
Rule & Admit & Pays & Value & Rev.\\\midrule
Capacity only & $A,B$ & $8,8$ & 10 & 16\\
EDF & $A,D$ & $0,0$ & 18 & 0\\
Feasible & $A,D$ & $9,6$ & 18 & 15\\
Credit $c=3$ & $A,E$ & $9,5$ & 16 & 14\\
Reserve $r=100$ & $A$ & 9 & 10 & 9\\
Both $r=100,c=5$ & $E$ & 5 & 6 & 5\\\bottomrule
\end{tabular}\end{center}
\noindent\textbf{Worked synthetic batch:} $T=400$ ms, each job 200 ms, values $(10,9,8,6)$ and deadlines $(200,200,400,400)$ for $A,B,D,E$; only $E$ is eligible. Capacity-only $B$ pays 8 for a late result and gains by bidding zero. Removing either queue or pose constraint permits $A,B$ and value 19. Credit 3 buys entrant access at a two-unit value cost. Across 300 paired six-robot batches (seed 206), mean values are 21.30, 38.86 and 43.06 for capacity-only, EDF and feasible allocation; Appendix A gives Monte Carlo intervals and 36,289 checks.

\textbf{Final self-test:} 88,200 joint bid/deadline deviations in 300 three-robot batches produce no positive gain. Appendix G explains the conditional result. In contrast, understating service time from 200 to 100 ms increases one operator's utility from 2 to 10. Verification must distinguish kinds of telemetry. The interface asks users to predict payment and timeliness separately before feedback \cite{hassidim}. Classroom questions motivate the test; automated interactions are not human-study evidence.

\balance
\section{Intellectual Merits}
RTC improves action continuity; Armory dynamically batches requests; Themis balances GPU efficiency and finish-time fairness. We do not reproduce or outperform those systems. Figure~\ref{fig:teaser} separates these foundations from our strategic counterexample, access trade-off and telemetry test. The critical-payment proof is established theory applied to this feasible family. Appendix G contains the three-paper comparison and expanded diagrams.

\section{Practical Impacts}
Operators could gain usable access and an appeal route for state errors. Academic evidence requires reproducibility; deployment requires calibrated latency and task success; public accountability requires transparent eligibility. The teaching artifact targets \textbf{SDG 4.4}: prediction accuracy before versus after a deadline explanation is the proposed indicator, not a measured benefit. Risks include exclusion through priority rules and confusing deadlines with safety. Publish access/value trade-offs and retain local safety control. The official-icon impact diagram is in Appendix G.

\subsection*{Future Research Roadmap}
Appendix G links the self-test to economics (audited duration reports), CS (matched Armory traces and solver overhead), and behavior (payment versus deadline explanations). Each study specifies data, a possible collaborator, a success criterion and a milestone. None is presented as completed robot or participant research.

\noindent\textbf{2056 aspiration.} Develop accountable access institutions whose incentives and physical benefits can be tested independently. The 2025 Economics Prize's innovation-and-growth foundation \cite{nobel2025} motivates studying who can use technological progress; our static model establishes no growth effect.

\noindent\textbf{Open artifacts.} \href{https://github.com/dku-comsci-econ206-Autumn2026/-PS2-FP9-ComputeCommons}{GitHub and embedded poster}; \href{https://huggingface.co/spaces/dku-comsci-econ206-2026/compute-commons}{Hugging Face}; \href{https://www.overleaf.com/project/6ac3a483f8fd778ea3fdc65f}{Overleaf}. Run \path{code/robot_model.py} and \path{code/deadline_test.py}. Release \texttt{final-fp9-2026-10-21}; tested source hashes and notebook are in the release record.
